% !PRINTED 273
\weberpage{273}
If we therefore set
\[
\begin{aligned}
A &= w^2 + w_1^2 + w_2^2,\\
B &= w^2w_1^2 + w^2w_2^2 + w_1^2w_2^2,\\
C &= w\,w_1\,w_2,
\end{aligned}
\]
we obtain
\[
\begin{array}{cccc}
\omega, & A, & B, & C, \\
-\dfrac{1}{\omega}, & A, & B, & C, \\
\omega+1, & e^{-\frac{\pi i}{12}(n-5)}A, & e^{-\frac{\pi i}{6}(n-5)}B, & C,
\end{array}
\]
so that two of these functions can be used just as $A$ and $B$ were used
above. For $n=5$ it appears, however, that in the expansions of $w$,
$w_1$, $w_2$ in powers of $q$ no negative powers occur, and that therefore
$A$, $B$, $C$, and hence also $w$, $w_1$, $w_2$ themselves, are constant.

The modular equation for $n=5$ can therefore be set up in each of the
three forms:
\weq{18}{
\begin{aligned}
u_1^2v_2^2 + u_2^2v_1^2 &= 2uv,\\
u_2^2v^2 - v_2^2u^2 &= 2u_1v_1,\\
u_1^2v^2 - v_1^2u^2 &= 2u_2v_2.
\end{aligned}
}

These equations can also be derived from the one given by \gesperrt{Jacobi}
(Fund. art. 30, gesammelte Werke, Bd. 1, S. 123).

We close these considerations by deriving, in the simplest cases, the
\gesperrt{Jacobi modular equations} from these irrational forms.

We set [cf. \S 54, (3)]:
\[
\begin{aligned}
\frac{u_2}{u} = x = \sqrt[4]{\varkappa}, &\qquad \frac{u_1}{u} = x' = \sqrt[4]{\varkappa'},\\[6pt]
\frac{v_2}{v} = y = \sqrt[4]{\lambda}, &\qquad \frac{v_1}{v} = y' = \sqrt[4]{\lambda'},
\end{aligned}
\]
% Source scan confirms the radicands \varkappa and \varkappa'.
and eliminate, by means of the relations
\weq{19}{
\begin{aligned}
x^8 + x'^8 &= 1, & xx' &= \frac{\sqrt{2}}{u^3}\\
y^8 + y'^8 &= 1, & yy' &= \frac{\sqrt{2}}{v^3}
\end{aligned}
}
% Source scan confirms the denominator v^3, parallel to u^3 above.
the quantities $u$ and $v$.
